\chapter[Sermon 13. The Lord Threatens Sore the Impenitent, as He That Renders to Every Man After His Works.]{The Lord Threatens Sore the Impenitent, as He That Renders to Every Man After His Works.}
\label{sermon:013}
\markright{Sermon 13. The Lord Threatens Sore the Impenitent, as He That ...}

And I gave her opportunity to repent of her fornication, but she did not repent. Behold, I will cast her into a bed, and those who commit fornication with her, into great adversity, unless they turn from their deeds. And I will put her children to death. And all congregations shall know that I am he who searches the minds and hearts, and I will give to each of you according to your works.

\index{Repentance}To the former errors and sins of Jezebel he adds another sin, nothing slight, namely, the abuse and even the contempt of God’s long-suffering. God does not immediately destroy those who are in error and sin, even the most grievous. But sinners are wont for the most part to abuse God’s long patience into an opportunity and pretext for sinning more openly, saying: “If God did so much abhor these offenses, he would have destroyed us long ago. But now he nourishes us benignly, therefore he does not greatly dislike it.” But this is an abuse of God’s long-suffering. For the Lord says at this present, “I have given Jezebel a time to repent, to leave her fornication, and to turn to the Lord.” However, she has not converted. Which thing the Lord takes in most evil part, that his grace should be truly despised and set at naught. Therefore, Paul to the Romans: “Do you despise the riches of God, his goodness, long-suffering, and lenity, not knowing that the goodness of God provokes you to repentance?” And so forth. If then the Lord has not suddenly oppressed us in our sins, let us not take that as license to sin, but let us rather amend. Peter says, “The Lord is patient toward us, while he will destroy none, but receive all to repentance.” 2 Peter 3. Certainly, Jezebel herself, when after the death of her husband Ahab, and the mortal fall of her son Ocozias, she did not amend, nor within the twelve years of her son Joram, wherein he is said to have reigned, did repeat her [illegible]—felt the wrath of God so much more grievous, for that it was long before it came.

And in the text following, the Lord Jesus in truth threatens the followers of Jezebel most severely—that is to say, the Cataphrigians or Montanists—unless they repent in time. For he opens again the gates of his grace to the penitent, recounting how he will punish the impenitent. By this, he truly seeks to drive them to repentance through threats. For in recounting the kinds or degrees of punishments, he also shows diverse kinds of those who are in error, and declares to each his judgment, which they may avoid through repentance. And it is thought that he rehearsed those kinds for this consideration, lest any man should perhaps think himself blameless and free, even if he is but a small partaker with Jezebel.

First, the Lord threatens Jezebel herself, that he will cast her into a bed. He speaks of the first authors of the evil and of the doctrine, upon whom he menaces to send a sickness. For the bed in many times in the scripture is taken for the very diseases with which those lying in bed are afflicted. And we Germans say that he is taken with a most grievous and deadly disease. And the Lord afflicts the archheretics with sickness of body and soul. In the meantime also he weakens the force of their terror, to the intent it might by little and little vanish away.

Secondly, he threatens great affliction to those who have dealings with Jezebel: that is to say, those who cleave to false doctrine, receive errors, delight in heresies, and seek to promote them. To these, he threatens most grievous afflictions, namely of body and soul, both in this present life and in the life to come. He seems to have said somewhat more than if he had merely recited certain kinds of punishment.

\index{Eusebius of Caesarea}Finally, he threatens death to the children born of this union and fornication, namely, illegitimate sons and bastards. These are chiefly the children of heretics, who stir up new heresies and revive those already condemned, weakened, and fading away. The Lord destroyed these with temporal and eternal death. And ecclesiastical history testifies that the Lord has indeed severely punished not only the heresy of the Cataphrigias but all heresies in general. Certain things concerning the Cataphrigias or Montanists are mentioned by Eusebius, book 1 of the \textit{Ecclesiastical History}, chapter 16.

The Lord seems to me here to have alluded to the old story of Jezebel and Ahab for them, as it were cast in a bed, from day to day, ever since they began to worship Baal; he vexed them with sickness and brought them low. The people who received the religion of Baal were put to much sorrow, evils, and afflictions. Finally, their children he brought to a shameful death. Their partakers also were slain, those who would have had Baal’s religion safe and sound, and even restored again. For after the death of Ahab his father, not many days after, Jehoaz the son of Ahab and Jezebel, bruised with an unhappy fall and cast in bed, died (2 Kings 1:4). And Jehoram, another son of Ahab and Jezebel, was slain, struck through with an arrow by Jehu (2 Kings 9). Athalia, the daughter of Ahab and Jezebel, the wife of Jehoram king of Judah, the son of Jehoshaphat, being ensnared by the oath of Jehoiada, fell down before the gates of the temple (2 Kings 11). And Jehoaz king of Judah, the son of Athalia and Jehoram, was slain also by the power of Jehu. And after were put to death by the same Jehu, seventy sons of Ahab. And all the priests of Baal were slain together in the temple, and before the altar of Baal, and not one of so great a number escaped. Yea, the temple, the idol, and the service of Baal were quite and clean overthrown. This old, marvelous, and wonderful history the Lord calls to memory, signifying that he lives yet a revenger and a punisher, who will neither overpass the just limit nor touch it out of time. For he adds, “And all congregations shall know,” etc.

\index{Ezekiel (Book of)}It is important and worthy of remembrance, and full of comfort, that in this recounting of punishments, he inserts a mention of repentance as though to say, let no one think that he must be destroyed and perish through a certain fatal necessity. For if anyone will repent, the gates of God’s grace are opened, his sins will be forgiven, and he will be received into favor and delivered from all those evils. And in this manner the Prophets have also taught, Jeremiah in chapter eighteen, and Ezekiel in chapter eighteen.

But since the punishment is not immediately executed upon the unrepentant, some will exclaim that God is asleep, that he sees or hears nothing. Therefore, the Lord himself answers them, saying: “And all congregations shall know…” For I will surely, at the last, execute my vengeance in due season. For then all men will learn that I neither sleep nor neglect my servants at any time, nor will I allow those who deserve evil of me and of my Church to escape unpunished. Furthermore, Christ testifies that he searches the depths and hearts of all men.

And he means that he knows all thoughts and devices of the heart, finally the appetite itself and all the desires of man, so that he can judge them truly, for nothing, however secret, is hidden from Christ. Therefore, he is very God. For it is the property of God, and belongs to him alone, to know the hearts of the children of men, as Solomon testifies in the third book of Kings. Christ therefore sees the secret and filthy works both of the Nicolaitans and all other beastly men, which Paul says are unworthy to come to light or to be expressly declared to men. Ephesians 5.

Christ does not only know all the thoughts of men, whatever they may be, but also gives to each man according to his own works. And so the Apostle Paul teaches, saying: “The just judgment of God will be revealed, which will reward each man according to his deeds—praise, honor, and immortality to those who continue in good works and seek eternal life, but indignation and wrath, tribulation, and anguish to those who are disobedient and disobey the truth, and follow iniquity.” 2 Romans. For works are the tests of faith and infidelity. And works, whether good or evil, are judged by God and by godly men, according as they proceed from faith or from infidelity. Therefore, whatever any of us sows, that he will also reap. For God is the most just rewarder of good and the revenger of evil. This sentence, as it is most true, is the foundation of the true and godly religion. Glory be to God.
